% Round 22 / geometry branch. Candidate; independent review required.
\section{A second complete-height family with maximal fourfold crossing depth}
\label{r22sq:section}

\paragraph{Scope.}
SU22 produces a second explicit family in the same ambient linear
system as DG22.  The local divisor arrangement now has genuine
quintuple points.  We compute the entire top and bottom weight
pieces and obtain exact fourth monodromy of rank $38864$.
Consequently these generic fibers are distinct from the DG22
family, where $N^3=0$.  Both families have full strong G and
complete height domains after every finite extension.  The
top weight coefficient is therefore not a Green obstruction.

\subsection{The actual five-divisor arrangement}

Work first over a finite field extension of $\mathbf F_5$.
Let $E:y^2=x^3-x$, $S=E^2$, $M=S\times S=E^4$, and put
\[
 A_0=\mathcal O_E(3[0])\boxtimes\mathcal O_E(3[0]),\qquad
 H=A_0\boxtimes A_0\quad\text{on }M.
 \tag{SQ22.1}\label{r22sq:H}
\]
Choose five smooth ample divisors $F_1,\ldots,F_5\in|H|$
meeting transversely, with all their intersections of one,
two or three distinct members geometrically integral and
with no fivefold intersection.  Every fourfold intersection
is a finite etale set.  After a further finite constant
extension split all of those sets into native points.
Such an arrangement exists: successive very ample Bertini
conditions, transversality to the finite list of old strata,
and avoidance of the old zero-dimensional intersections
define a nonempty open over the algebraic closure and hence
over a finite constant field.  The stated geometric
integrality is part of these Bertini conditions.

Each fourfold intersection consists of exactly
\[
                 c:=\deg_M(H^4)=4!\,3^4=1944
 \tag{SQ22.2}\label{r22sq:count}
\]
native points.  Indeed $c_1(H)$ is the sum of the four
degree-three point classes of the elliptic factors; only
their product survives in the fourth power.  The transversality
makes every local multiplicity one.

Choose five distinct native points of $\mathbf P^1$.
The ten factor divisors on $M\times\mathbf P^1$ are
\[
 F_i\times\mathbf P^1\quad(1\le i\le5),\qquad
 M\times\{p_j\}\quad(1\le j\le5),
 \tag{SQ22.3}\label{r22sq:factors}
\]
ordered with the first five before the point factors.  Choose
$Q\in H^0(M\times\mathbf P^1,H^5\boxtimes\mathcal O(5))$
with all the transverse-cut hypotheses of SU22, in particular
avoiding the zero-dimensional fivefold strata.  The same
simultaneous Bertini argument supplies it after another finite
constant extension.  The positive-dimensional center cuts are
geometrically integral componentwise.  Zero-dimensional center
cuts are only required to be finite etale.

Lift this datum to a good mixed-characteristic trait and
assume the resulting generic divisor
\[
          X=V\left((\prod_{i=1}^5 f_i)(\prod_{j=1}^5 l_j)
                                      +\pi Q\right)
 \tag{SQ22.4}\label{r22sq:smoothing}
\]
is smooth and geometrically integral.  SU22 supplies full G
in every degree and after every finite valued extension.
The actual ordered resolution keeps the ten components,
and each stratum is the codimension-two smooth blowup
prescribed there.  The quadruple strata are unchanged curves:
either a disjoint union of $c$ copies of $\mathbf P^1$
(four $F$ factors), or one smooth geometrically integral
curve $F_i\cap F_j\cap F_k$ (three $F$ factors and a point).
They have no center of codimension two.  The quintuple
strata are exactly the $c$ native points for each choice
of four $F$ factors and one point factor.  Thus there are
$25c$ such points and no sixfold strata.

\subsection{The whole extreme coefficient complex}

Let $A$ denote the finite incidence complex of the five
divisors on $M$.  It has five vertices, ten edges, ten
triangular faces, and $c$ distinct tetrahedral cells for
each of the five four-element subsets of $\{1,\ldots,5\}$.
Every such cell has the ordinary oriented boundary of its
four-element subset.  These repeated cells record the
distinct native intersection points.  This is an incidence
or delta-complex; it is not incorrectly treated as a simple
abstract simplicial complex with only five tetrahedra.

\begin{lemma}[Exact top incidence dimension]
\label{r22sq:incidence}
The reduced rational homology of $A$ is zero below degree
three, and
\[
                         \dim H_3(A,\mathbf Q)=5c-4.
 \tag{SQ22.5}\label{r22sq:H3}
\]
The dual incidence complex $\Delta$ of the actual model
\eqref{r22sq:smoothing} is the join of $A$ with five discrete
vertices.  Consequently its reduced rational cohomology is
concentrated in degree four and
\[
                     \dim H^4(\Delta,\mathbf Q)=4(5c-4)=38864.
 \tag{SQ22.6}\label{r22sq:H4}
\]
\end{lemma}

\begin{proof}
Through dimension two the complex $A$ is the full
two-skeleton of the simplex on five vertices.  Its
tetrahedral boundary map has $5c$ columns, consisting of
$c$ copies of each of the five usual columns for the
boundary of the four-simplex.  That usual boundary map
has rank four: the single relation is the oriented
boundary of the four-simplex itself, and the full simplex
chain complex is exact.  Repeating columns changes neither
this rank nor its image.  The image is exactly the kernel
of the triangular boundary map, and the lower reduced
homology is zero by the same simplex exactness.  There
are no four-cells in $A$.  Hence its third homology has
dimension $5c-4$.

For a stratum indexed only by $F$ factors, its old form
is $F_I\times\mathbf P^1$; for a stratum also containing
one point factor, it is $F_I$.  There is no stratum with
two point factors.  The geometric connected components
and every incidence map are therefore precisely the join
data.  The actual smooth blowups do not change connected
components, and the resolution adds no components or
incidences; its local charts and the explicit SU22 stratum
formula give the same join.  In particular the $c$ labels
of each fourfold $F$ intersection remain identified under
all five point-face maps.

The join is a union of five cones on $A$ with the same
common base.  Its reduced chain complex is equivalently
the suspended reduced chain complex of $A$ tensor the
four-dimensional reduced degree-zero complex of five
points.  Thus
$\widetilde H_j(\Delta)=\widetilde H_{j-1}(A)^{\oplus4}$.
This proves the asserted entire reduced homology and
cohomology statement.  As a count check its cell numbers
are
\[
                (c_0,c_1,c_2,c_3,c_4)
                      =(10,35,60,5c+50,25c),
\]
whose Euler characteristic is $1+4(5c-4)$.
\end{proof}

\begin{theorem}[Exact fourth monodromy, SQ22-local]
\label{r22sq:N4}
For every supplied smoothing \eqref{r22sq:smoothing}, the
actual generic cohomology satisfies
\[
                 \operatorname{rank}(N^4|H^4)=38864,\qquad N^5=0.
 \tag{SQ22.7}\label{r22sq:monodromy}
\]
The minimum strict semistable dominance depth, taken over
all finite valued extensions and all smooth projective
dominant witnesses, is exactly four.  The lower bound
holds over each fixed finite valued extension.
\end{theorem}

\begin{proof}
Use the Saito convention explicitly displayed in DW22.
There are no sixfold strata.  The whole lower extreme is
the last cohomology of the $H^0$ incidence row:
\[
 E_2^{4,0}=\operatorname{coker}
   \{\delta_3:H^0(D^{[4]})\longrightarrow H^0(D^{[5]})\}
                    =H^4(\Delta,\mathbf Q_\ell).
 \tag{SQ22.8}\label{r22sq:lower}
\]
All $H^0$ groups here are indexed by actual geometric
connected components, which are native in the chosen
preparation.  The entire upper extreme is
\[
 E_2^{-4,8}=\ker\{H^0(D^{[5]})(-4)
                         \longrightarrow H^2(D^{[4]})(-3)\}.
 \tag{SQ22.9}\label{r22sq:upper}
\]
There is no incoming upper quotient.  Every quadruple
component is one of the smooth geometrically integral
curves listed above.  Its $H^2$ has the trace-normalized
point generator.  Every quintuple point is native and
has degree one on that curve.  The upper Gysin map is
therefore exactly $\delta_3^t$ under these generators,
with the common incidence sign convention.  No ambient
Gysin was cancelled in reaching this assertion.

The common-stratum rule for $N$ makes $N^4$ the identity
on the quintuple $H^0$ summands, between the two extreme
positions.  It induces
\[
 \ker\delta_3^t\longrightarrow
                    \mathbf Q_\ell^{25c}/\operatorname{im}\delta_3.
 \tag{SQ22.10}\label{r22sq:harmonic}
\]
This map is an isomorphism.  Over $\mathbf Q$, the
ordinary positive dot product on the native-point basis
gives
$\mathbf Q^{25c}=\operatorname{im}\delta_3\oplus
\ker\delta_3^t$.  Tensoring with $\mathbf Q_\ell$ preserves
that direct sum.  Thus its dimension is exactly the one
computed in \eqref{r22sq:H4}.  The argument does not assign
a positive ordering to an $\ell$-adic form.

Finite-field purity gives degeneration at $E_2$ exactly
as in DW22.  The abutment on $H^4$ has weights in $[0,8]$;
$N^4$ kills $W_7$ and lands in $W_0(-4)$, while $W_{-1}=0$.
Hence the actual operator's rank is precisely the rank
of \eqref{r22sq:harmonic}.  A fifth application is zero.
No general weight--monodromy theorem is needed.

Finite extension changes the normalized monodromy by a
nonzero ramification scalar, or leaves it unchanged with
the transported tame character; its rank is invariant.
A smooth projective dominant witness of depth at most
three would have $N^4=0$ on $H^4$, because its spectral
sequence has no columns outside $[-3,3]$.  Its pullback
on $H^4$ from $X$ is injective by the positive-degree
ample trace calculation DW22.12.  This contradicts
\eqref{r22sq:monodromy}.  The actual model itself has
depth four and gives a witness over the original field.
The stated minimum is therefore four.
\end{proof}

\subsection{A separate all-place arithmetic construction}

\begin{theorem}[Quintuple complete-height family, SQ22]
\label{r22sq:global}
For the everywhere-good global data of DG22,
$\mathcal Y=\mathcal E^4\times\mathbf P^1$ and
$\mathcal L_0=(\mathcal A_0\boxtimes\mathcal A_0)
\boxtimes\mathcal O(1)$, there is one fixed finite field
extension $F/F_0$, unramified over $5$, and a Zariski-dense
family in $|\mathcal L_0^5|_{Y_F}$ such that:
\begin{enumerate}
\item every specified geometrically homologically trivial
rational cycle, in every codimension, has full strong G
at every finite place after every finite $F'/F$, including
every new cycle over the corresponding completion;
\item PH16 defines the complete complementary BB height
pairings on
\[
 \operatorname{CH}^r(X_{F'})_{\mathbf Q,\mathrm{geom\,hom}}
 \times\operatorname{CH}^{5-r}(X_{F'})_{\mathbf Q,\mathrm{geom\,hom}},
                         \qquad1\le r\le4,
 \tag{SQ22.11}\label{r22sq:height}
\]
and GA17 gives admissible regular-model Chow lifts on
these full proved domains;
\item at every distinguished place above $5$ its exact
fourth monodromy has rank $38864$ and its minimum strict
semistable dominance depth is four in the scope of
Theorem~\ref{r22sq:N4};
\item no member is geometrically isomorphic to a member
of DG22, even after a common finite extension.  It is
also distinct from the named SG18/MG21 families and
projective-space complete-intersection fourfolds by
$b_1=8,b_2=29$.
\end{enumerate}
The family is dense over one fixed $F$; the whole linear
system and arbitrary varieties are not claimed to satisfy G.
\end{theorem}

\begin{proof}
We specify the local replacement and recheck the arithmetic
steps, so that a theorem about one completion is not
mistaken for a complete height construction.

Use DG22's full-section coefficient module $V_5$ for this
same ambient.  Choose a smooth generic integral reference
$H_0$ and the finite rational-prime set $S_0$ of all its
bad reductions other than $5$, together with $2$.
At every place over every prime in $S_0$ choose a smooth
residue replacement $G_v$.  At every place over $5$ choose
the actual five-divisor arrangement above, split all its
fourfold intersection points, choose the five projective
points, and choose $Q_v$ with all SU22 hypotheses.  These
are finitely many conditions over finite residue
extensions.  A monic irreducible CRT polynomial of one
degree divisible by all their degrees gives a single
fixed extension $F/F_0$ unramified and inert at all the
selected places, exactly as proved in DG22.

Positive elliptic cohomology on $M=E^4$ makes sections
of $H=\mathcal A_0\boxtimes\mathcal A_0$ commute with
base change.  CRT therefore lifts each of the five
divisor sections and each of the five point forms.
Their actual product is
\[
 P=(\prod_{i=1}^5 f_i)(\prod_{j=1}^5l_j)
                 \in H^0(\mathcal Y,\mathcal L_0^5).
 \tag{SQ22.12}\label{r22sq:globalP}
\]
Relative to $H$ and $\mathcal O(1)$ it has bidegree
$(5,5)$, hence has a degree-five ambient polynomial
representative by DG22's coefficient lemma.  The
individual $f_i$ and $l_j$ are not incorrectly treated
as degree-one Segre forms.  Lift the $Q_v$ and the
smooth replacements in the same free polynomial module.

Put $B=\prod_{p\in S_0}p$, $M_0=25B$, and require
\[
 H_1\equiv P+5Q\pmod{25\mathcal O_F},\qquad
 H_1\equiv G_v\pmod{\mathfrak p_v}\quad(v\text{ above }S_0).
 \tag{SQ22.13}\label{r22sq:H1}
\]
The solution lattice is dense.  Intersect it with the
same nonempty fixed-reference Lefschetz pencil open of
DG22; the ambient, degree and jet surjectivity are
unchanged.  After fixing $H_1$, spread this pencil and
let $T_0$ contain all exceptional rational primes.
Choose infinitely many integers $a$ with
$a\equiv1\pmod{25B}$ and $a\equiv0\pmod p$ for
$p\in T_0\setminus(S_0\cup\{5\})$, omitting the
finitely many singular generic parameters, and set
\[
                            X_a=V(M_0H_0+aH_1).
 \tag{SQ22.14}\label{r22sq:Xa}
\]

At every place over $5$, dividing by the unit $a$
gives $P+5Q_a$ with
$Q_a=Q+5R_1+5B a^{-1}H_0$ and $\bar Q_a=\bar Q$.
It is precisely the supplied SU22/SQ22 model, since
$5$ is a native uniformizer.  At places over $S_0$ the
special section is $G_v$.  At any other exceptional
place it is a nonzero multiple of the smooth reference
$H_0$.  At every remaining place the spread pencil
has a smooth or one-node special divisor.  All
characteristic-two places were included in $S_0$.
Thus there is no omitted finite place.

The relative Cartier and geometric integrality checks
of DG22 apply literally: the ambient and line bundle
are identical, each defining special section is
nonzero, and $H^0(\mathcal L_0^{-5})=
H^1(\mathcal L_0^{-5})=0$ on every geometric fiber.
Smooth or isolated-node fourfold divisors are connected
and normal and hence integral.  The flat closures are
actual projective integral models.

At $5$, SU22 applies to every new cycle after every
finite valued extension; at other places the frozen
good-reduction and ON17 all-finite-extension theorems
apply.  Consequently for an arbitrary finite $F'/F$
all local cycle groups on all completions satisfy G
anew.  The full exact supports, all ordinary actual
tests, full field degrees and extending original
absolute values are those of the frozen interfaces.
PH16 and its closed archimedean input now give
\eqref{r22sq:height}; GA17 gives the admissible lifts.

The fixed field and residue lattice were chosen before
varying $H_1$.  For every allowed direction there are
infinitely many allowed $a$.  The affine-line vanishing
argument in DG22 therefore proves density over this
one fixed $F$ in the full coefficient space and then
in the full section space.  Changing the finite
exceptional set with the direction causes no problem.

Finally Theorem~\ref{r22sq:N4} gives the exact local
rank.  DG22 has $N^3=0$ and hence $N^4=0$ at every
place over $5$.  A geometric isomorphism descends to
a finite extension and would identify these Galois
representations over a common finite valued extension;
nonzero versus zero $N^4$ contradicts that possibility.
Weak Lefschetz gives $b_1=8,b_2=29$ exactly as in
DG22 and makes the other named distinctions.
The proof establishes a genuinely different family
of generic fibers, not merely another model of DG22.
\end{proof}

\paragraph{Contribution and dependencies.}
The complete numerical fourfold gate is proved here
on actual models with maximal possible crossing depth.
The full final Green test space is unchanged.  The
new local input is SU22; the exact top-incidence and
monodromy calculation is given above.  The all-place
arithmetic construction uses the explicitly verified
DG22 coefficient and pencil mechanism with a different
local residue target.  The retained frozen dependencies
are FP21, SP19, PU18, CD17/CL17, MZ20, RI18, VT17,
ON17, PH16 and GA17.  Saito's spectral sequence and
Deligne's purity are used with the exact conventions
and primary-source checks of DW22.  General G and
arbitrary-variety complete heights remain open;
no literature-priority claim is made.
